\chapter{Giới thiệu}
\label{ch:intro}

Chương này trình bày bối cảnh và động lực của khóa luận, phân tích hai hạn chế cốt lõi của mô
hình ngôn ngữ lớn khi giải quyết các câu hỏi có ràng buộc về thời gian, chỉ ra khoảng trống
nghiên cứu trong bối cảnh tiếng Việt, từ đó phát biểu mục tiêu, câu hỏi nghiên cứu và các đóng
góp chính. Cuối chương là cấu trúc tổng thể của khóa luận.

\section{Bối cảnh và động lực}
\label{sec:motivation}

Trong những năm gần đây, các mô hình ngôn ngữ lớn (\textit{Large Language Models} -- LLM) như
GPT, LLaMA hay Qwen đã tạo nên bước tiến vượt bậc cho hàng loạt tác vụ xử lý ngôn ngữ tự nhiên,
trong đó có bài toán hỏi đáp tri thức (\textit{Knowledge Question Answering})~\cite{zhao2023llmsurvey}.
Nhờ được huấn luyện trên kho ngữ liệu khổng lồ, LLM có khả năng trả lời trôi chảy nhiều câu hỏi
thông thường mà không cần truy cập nguồn tri thức bên ngoài. Tuy nhiên, khi câu hỏi đòi hỏi
\textbf{suy luận theo thời gian} -- tức là phải xác định, so sánh và lập luận trên các sự kiện
gắn với những mốc thời gian khác nhau -- thì hiệu năng của LLM suy giảm rõ rệt và mô hình thường
đưa ra câu trả lời thiếu chính xác hoặc gây hiểu nhầm~\cite{chen2024ari}.

Nguyên nhân sâu xa nằm ở bản chất của tri thức thực tế: tri thức không tĩnh tại mà
\textbf{tiến hóa theo thời gian}. Cùng một câu hỏi nhưng đáp án thay đổi tùy thời điểm. Ví dụ,
quốc gia đăng cai Thế vận hội Mùa đông năm 2018 là Hàn Quốc, nhưng đến năm 2022 lại là Trung
Quốc. Một thực thể có thể giữ nhiều vai trò ở những giai đoạn khác nhau, và mối quan hệ giữa các
thực thể luôn gắn với một khoảng hiệu lực nhất định. Để trả lời đúng, hệ thống không những phải
\emph{truy xuất} đúng sự kiện mà còn phải \emph{xử lý} chúng theo trình tự logic chặt chẽ.

Hãy xét ví dụ trong Hình~\ref{fig:temporal-example}: với câu hỏi \textit{``Tổng thống Hoa Kỳ khi
Chiến tranh Việt Nam kết thúc là ai?''}, hệ thống cần thực hiện nhiều bước suy luận: (1) xác định
thời điểm Chiến tranh Việt Nam kết thúc (năm 1975); (2) truy vấn ai là Tổng thống Hoa Kỳ tại
thời điểm đó; (3) trả về thực thể tương ứng làm đáp án. Những câu hỏi dạng này yêu cầu kết hợp
đồng thời thông tin về thực thể, quan hệ và mốc thời gian -- vốn được biểu diễn tự nhiên trong
một \textbf{đồ thị tri thức thời gian} (\textit{Temporal Knowledge Graph} -- TKG).

\begin{figure}[ht]
  \centering
  % \includegraphics[width=0.85\textwidth]{ch1_temporal_example.pdf}
  \fbox{\parbox{0.85\textwidth}{\centering\vspace{3mm}
  \textbf{Câu hỏi:} ``Tổng thống Hoa Kỳ khi Chiến tranh Việt Nam kết thúc là ai?''\\[2mm]
  \footnotesize
  Bước 1: Chiến~tranh~Việt~Nam $\xrightarrow{\text{kết thúc}}$ 1975 \\
  Bước 2: ?~$\xrightarrow{\text{là tổng thống Hoa Kỳ}}$ (1974--1977) \\
  Bước 3: $\Rightarrow$ \textbf{Gerald Ford}
  \vspace{3mm}}}
  \caption{Một câu hỏi có ràng buộc về thời gian đòi hỏi chuỗi suy luận nhiều bước trên đồ thị
  tri thức thời gian. \todo{thay bằng hình minh họa TKG đẹp hơn}}
  \label{fig:temporal-example}
\end{figure}

Bài toán hỏi đáp trên đồ thị tri thức thời gian (\textit{Temporal Knowledge Graph Question
Answering} -- TKGQA) vì vậy đã trở thành một hướng nghiên cứu quan trọng~\cite{saxena2021cronkgqa,
chen2023multitq}. Các phương pháp truyền thống dựa trên luật được thiết kế thủ công hoặc trên
biểu diễn nhúng học được, tuy hiệu quả với câu hỏi đơn giản nhưng khó mở rộng và thường thất bại
với các câu hỏi đòi hỏi suy luận thời gian tinh vi~\cite{chen2022subgraph}. Sự xuất hiện của LLM
mở ra cơ hội tận dụng năng lực suy luận linh hoạt của mô hình, song như đã nêu, việc dùng LLM một
cách trực tiếp lại vấp phải những hạn chế căn bản sẽ phân tích ở mục tiếp theo.

\section{Hai hạn chế cốt lõi của mô hình ngôn ngữ lớn}
\label{sec:limitations}

Qua phân tích trực giác và thực nghiệm, có thể quy các khó khăn của LLM trong suy luận thời gian
về hai nguyên nhân chính~\cite{chen2024ari}.

\paragraph{Thiếu tri thức thời gian.} LLM thu nhận tri thức thông qua quá trình tiền huấn luyện
trên một kho dữ liệu cố định. Sau khi huấn luyện, các tham số của mô hình bị ``đóng băng'', khiến
mô hình không thể cập nhật những tri thức mới hoặc tri thức biến đổi sau thời điểm huấn luyện. Hệ
quả là LLM dễ sinh ra hiện tượng \textit{ảo giác} (\textit{hallucination}) -- bịa ra thông tin
không có thật -- khi gặp các sự kiện nằm ngoài phạm vi dữ liệu huấn luyện, hoặc khi tri thức đã
thay đổi so với thời điểm mô hình được học.

\paragraph{Thiếu suy luận thời gian phức tạp.} Do cơ chế sinh đầu ra dựa trên xác suất cực đại
theo từng bước, LLM gặp khó khăn khi phải thực hiện chuỗi suy luận nhiều bước liên kết chặt chẽ
với nhau. Trong một câu hỏi thời gian phức tạp, chỉ cần một sai sót ở khâu truy xuất tri thức hay
một nhầm lẫn ở khâu lọc theo thời gian cũng đủ khiến cả chuỗi lập luận dẫn tới kết luận sai. Các
lỗi nhỏ tích lũy dần qua từng bước sinh khiến độ tin cậy của kết quả cuối cùng giảm mạnh.

Các hướng tiếp cận hiện có thường chỉ giải quyết được một phần. Việc bổ sung tri thức ngoài theo
kiểu sinh tăng cường bằng truy hồi (\textit{Retrieval-Augmented Generation} --
RAG)~\cite{lewis2020rag, baek2023kaping} giúp làm giàu thông tin cho LLM, nhưng độ chính xác của
khâu truy hồi cùng giới hạn độ dài ngữ cảnh lại dễ đưa thêm nhiễu và làm thiếu manh mối suy luận.
Việc cung cấp ví dụ mẫu (\textit{exemplar}) để dẫn dắt LLM~\cite{wei2022cot} thì phụ thuộc vào
chất lượng ví dụ và tốn nhiều công sức xây dựng, đồng thời khó bao quát sự đa dạng của các câu hỏi
thực tế. Như minh họa trong Hình~\ref{fig:three-levels}, điểm chung của các hướng này là chưa cung
cấp được \emph{hướng dẫn ở mức phương pháp luận} cho quá trình suy luận thời gian đang diễn ra, và
dễ bị nhiễu bởi tri thức cụ thể không liên quan.

\begin{figure}[ht]
  \centering
  % \includegraphics[width=0.9\textwidth]{ch1_three_levels.pdf}
  \fbox{\parbox{0.9\textwidth}{\centering\vspace{2mm}\footnotesize
  \textbf{Tri thức cụ thể} (Specific) $\rightarrow$ truy hồi sự kiện, phạm vi áp dụng hẹp\\
  \textbf{Ví dụ mẫu} (Exemplar) $\rightarrow$ học trong ngữ cảnh, phạm vi áp dụng vừa\\
  \textbf{Trừu tượng} (Abstract) $\rightarrow$ hướng dẫn phương pháp luận, phạm vi áp dụng rộng
  \vspace{2mm}}}
  \caption{Ba mức tích hợp thông tin với mô hình ngôn ngữ lớn. Tri thức càng trừu tượng và tinh
  lọc thì phạm vi áp dụng càng rộng (phỏng theo~\cite{chen2024ari}). \todo{vẽ lại bằng TikZ}}
  \label{fig:three-levels}
\end{figure}

Phương pháp \textbf{Suy luận Trừu tượng Quy nạp} (\textit{Abstract Reasoning Induction} -- ARI)
do Chen và cộng sự đề xuất~\cite{chen2024ari} hướng tới khắc phục đồng thời cả hai hạn chế trên.
Lấy cảm hứng từ \textit{thuyết kiến tạo} (\textit{constructivism}) trong giáo
dục~\cite{savery1995constructivism}, ARI chia quá trình suy luận thời gian thành hai pha tách
bạch: pha \textit{không phụ thuộc tri thức} (\textit{knowledge-agnostic}), nơi LLM chỉ đưa ra
quyết định chiến lược ở mức trừu tượng; và pha \textit{dựa trên tri thức}
(\textit{knowledge-based}), nơi các thao tác cụ thể được thực thi trên đồ thị tri thức. Cách tách
biệt này vừa cung cấp chỗ dựa tri thức xác thực cho LLM, vừa giảm thiểu nhiễu từ dữ liệu cụ thể.
Đồng thời, ARI cho phép mô hình \emph{chủ động học} các chiến lược suy luận trừu tượng từ những
lần giải trước đó -- cả đúng lẫn sai -- để tái sử dụng khi gặp câu hỏi tương tự. Đây chính là nền
tảng phương pháp mà khóa luận kế thừa và phát triển.

\section{Khoảng trống nghiên cứu trong bối cảnh tiếng Việt}
\label{sec:gap}

Mặc dù TKGQA đã được nghiên cứu khá sâu cho tiếng Anh, bối cảnh tiếng Việt vẫn còn nhiều khoảng
trống đáng kể.

\textbf{Thứ nhất, thiếu bộ dữ liệu TKGQA tiếng Việt.} Các bộ dữ liệu chuẩn như
CronQuestions~\cite{saxena2021cronkgqa} hay MultiTQ~\cite{chen2023multitq} đều được xây dựng cho
tiếng Anh, dựa trên các nguồn như Wikidata hoặc ICEWS với nhãn tiếng Anh. Cho đến nay chưa có một
bộ dữ liệu hỏi đáp trên đồ thị tri thức thời gian nào dành riêng cho tiếng Việt, khiến việc huấn
luyện, đánh giá và so sánh các phương pháp cho tiếng Việt gặp trở ngại lớn.

\textbf{Thứ hai, phụ thuộc vào mô hình ngôn ngữ thương mại.} Phần lớn các hệ thống TKGQA dựa trên
LLM hiện nay vận hành thông qua các giao diện lập trình ứng dụng (API) thương mại như
\texttt{gpt-3.5-turbo}~\cite{chen2024ari}. Cách làm này phát sinh chi phí, phụ thuộc kết nối mạng
và đặt ra lo ngại về quyền riêng tư dữ liệu, đồng thời gây khó khăn cho việc triển khai
\emph{cục bộ} trong những môi trường hạn chế tài nguyên hoặc yêu cầu bảo mật.

\textbf{Thứ ba, chưa khai thác nhận dạng thực thể tiếng Việt cho liên kết thực thể.} Để áp dụng
TKGQA cho câu hỏi mở do người dùng nhập, hệ thống cần \textit{liên kết thực thể} (\textit{entity
linking}) -- ánh xạ cụm từ trong câu hỏi về đúng nút trong đồ thị tri thức. Các mô hình nhận dạng
thực thể có tên (\textit{Named Entity Recognition} -- NER) tiếng Việt đã đạt chất lượng tốt nhưng
chưa được tận dụng cho bài toán TKGQA tiếng Việt.

Những khoảng trống trên vừa là thách thức, vừa là động lực để khóa luận xây dựng một hệ thống
TKGQA hoàn chỉnh cho tiếng Việt: từ việc kiến tạo dữ liệu, thích nghi phương pháp suy luận, đến
cải tiến cơ chế lọc hành động.

\section{Mục tiêu và câu hỏi nghiên cứu}
\label{sec:objectives}

\textbf{Mục tiêu tổng quát} của khóa luận là nghiên cứu và áp dụng phương pháp Suy luận Trừu
tượng Quy nạp nhằm nâng cao năng lực trả lời các câu hỏi có ràng buộc về thời gian cho tiếng Việt,
đồng thời xây dựng một hệ thống thử nghiệm hoàn chỉnh và một bộ dữ liệu đánh giá đi kèm.

Để cụ thể hóa, khóa luận đặt ra ba câu hỏi nghiên cứu:

\begin{itemize}
  \item \textbf{CHNC1:} Phương pháp ARI có cải thiện độ chính xác so với việc dùng mô hình ngôn
  ngữ lớn suy luận trực tiếp trên bài toán TKGQA tiếng Việt hay không, và cải thiện ở mức nào?
  \item \textbf{CHNC2:} Mỗi thành phần trong khung ARI (hướng dẫn trừu tượng, phân cụm lịch sử,
  lọc hành động, học từ ví dụ sai) đóng góp ra sao vào hiệu năng tổng thể?
  \item \textbf{CHNC3:} Bộ lọc hành động \emph{nhận biết ngữ cảnh} dựa trên cross-encoder có khắc
  phục được hạn chế của bộ lọc dựa trên độ tương đồng cosine truyền thống hay không?
\end{itemize}

\section{Đóng góp của khóa luận}
\label{sec:contributions}

Khóa luận có ba đóng góp chính, tương ứng với ba câu hỏi nghiên cứu nêu trên.

\begin{enumerate}
  \item \textbf{Xây dựng CronQ-VN -- bộ dữ liệu TKGQA tiếng Việt.} Khóa luận thiết kế một quy
  trình trích xuất đồ thị tri thức thời gian từ Wikidata và sinh câu hỏi tự động, tạo ra bộ dữ
  liệu CronQ-VN với khoảng 238 nghìn sự kiện, 37,5 nghìn thực thể và 271 quan hệ. Trên cơ sở đó,
  21 quan hệ giàu dữ liệu thuộc nhiều lĩnh vực được chọn để sinh năm loại câu hỏi thời gian, với
  đáp án chuẩn được truy vấn trực tiếp từ đồ thị. Phần nhãn còn thiếu tiếng Việt được bổ sung bằng
  mô hình ngôn ngữ lớn.

  \item \textbf{Thích nghi khung ARI cho tiếng Việt với mô hình ngôn ngữ cục bộ.} Khóa luận cài
  đặt lại toàn bộ khung ARI để vận hành hoàn toàn bằng LLM chạy cục bộ thông qua
  Ollama~\cite{ollama} với mô hình Qwen2.5~\cite{qwen2025}, thay cho các API thương mại. Hệ thống
  được bổ sung cơ chế liên kết thực thể cho câu hỏi mở, sử dụng mô hình NER tiếng Việt
  (\texttt{NlpHUST/ner-vietnamese-electra-base}) kết hợp tìm kiếm tương đồng trong cơ sở dữ liệu
  vector.

  \item \textbf{Đề xuất bộ lọc hành động dựa trên cross-encoder.} Khóa luận chỉ ra hạn chế của cơ
  chế lọc hành động bằng độ tương đồng cosine trong ARI gốc -- vốn so khớp hành động với câu hỏi
  gốc mà không nhận biết bước suy luận hiện tại -- và đề xuất thay thế bằng một bộ lọc nhận biết
  ngữ cảnh dựa trên cross-encoder~\cite{reimers2019sbert}, mã hóa đồng thời hành động và lịch sử
  suy luận để chấm điểm mức phù hợp.
\end{enumerate}

Về phạm vi, khóa luận tập trung vào các câu hỏi yêu cầu suy luận nhiều bước trên đồ thị con của
đồ thị tri thức thời gian, với dữ liệu chính là bộ CronQ-VN tự xây; các bộ dữ liệu tiếng Anh như
CronQuestions và MultiTQ được dùng để tham chiếu và đối sánh phương pháp.

\section{Cấu trúc khóa luận}
\label{sec:structure}

Phần còn lại của khóa luận được tổ chức như sau:

\begin{itemize}
  \item \textbf{Chương~\ref{ch:background} -- Cơ sở lý thuyết} trình bày các kiến thức nền: đồ thị
  tri thức thời gian, bài toán TKGQA, mô hình ngôn ngữ lớn và học trong ngữ cảnh, tăng cường tri
  thức, suy luận có cấu trúc, bộ nhớ và thuyết kiến tạo, phân cụm K-means, cross-encoder và nhận
  dạng thực thể tiếng Việt.
  \item \textbf{Chương~\ref{ch:related} -- Các công trình liên quan} khảo sát các mô hình TKGQA,
  các hướng dùng LLM kết hợp tri thức ngoài và bộ nhớ, các bộ dữ liệu TKGQA, và định vị đề tài
  trong bức tranh chung.
  \item \textbf{Chương~\ref{ch:method} -- Phương pháp đề xuất} mô tả chi tiết kiến trúc hệ thống
  hai pha, quy trình xây dựng CronQ-VN, cơ chế thích nghi tiếng Việt, và bộ lọc cross-encoder.
  \item \textbf{Chương~\ref{ch:experiment} -- Thực nghiệm và Kết quả} trình bày thiết lập thực
  nghiệm, các phương pháp so sánh, độ đo và phân tích kết quả theo ba câu hỏi nghiên cứu.
  \item \textbf{Chương~\ref{ch:conclusion} -- Kết luận và Hướng phát triển} tổng kết đóng góp, nêu
  ý nghĩa, hạn chế và các hướng mở rộng trong tương lai.
\end{itemize}
